%% ============================================================
%% EK F — Kaynakça ve İleri Okuma
%% Olasılık Teorisi ve Matematiksel İstatistik
%% ============================================================
\chapter{Kaynakça ve İleri Okuma}
\label{app:kaynaklar}

\section{Temel Kaynaklar}

Bu kitabın yazımında doğrudan yararlanılan başvurular:

\begin{description}
  \item[Akdeniz (2013)] Fikri Akdeniz. \textit{Olasılık ve İstatistik.} 18. Baskı. Çukurova Üniversitesi, Adana. — Türkçe bağlam için temel referans; örnek ve alıştırmalar.

  \item[Billingsley (1995)] Patrick Billingsley. \textit{Probability and Measure.} 3rd ed. Wiley. — Ölçüm teorisi ve olasılığın kapsamlı birleşimi; Böl.~0--2 için.

  \item[Casella-Berger (2002)] George Casella, Roger Berger. \textit{Statistical Inference.} 2nd ed. Duxbury. — İstatistik Part II'nin ana kaynağı; yeterlilik, UMP, regresyon.

  \item[Durrett (2010)] Rick Durrett. \textit{Probability: Theory and Examples.} 4th ed. Cambridge. — Martingaller, GBSK, MLT için birincil kaynak.

  \item[Feller Vol.I (1968)] William Feller. \textit{An Introduction to Probability Theory and Its Applications.} Vol.~I, 3rd ed. Wiley. — Kesikli olasılık ve limit teoremleri; sezgi için eşsiz.

  \item[Feller Vol.II (1971)] William Feller. \textit{An Introduction to Probability Theory and Its Applications.} Vol.~II, 2nd ed. Wiley. — Sürekli dağılımlar, Laplace dönüşümü, yenileme.

  \item[Gnedenko (1962)] Boris Gnedenko. \textit{The Theory of Probability.} Chelsea. — Klasik Rus ekolü; limit teoremlerinin tarihsel kökeni.

  \item[Hogg-Craig (1995)] Robert Hogg, Allen Craig. \textit{Introduction to Mathematical Statistics.} 5th ed. Prentice Hall. — İstatistik Part II için tamamlayıcı kaynak.

  \item[Klenke (2014)] Achim Klenke. \textit{Probability Theory: A Comprehensive Course.} 2nd ed. Springer. — Modern ölçüm teorisi yaklaşımı; Böl.~0 için ek referans.

  \item[Kolmogorov (1950)] Andrei Kolmogorov. \textit{Foundations of the Theory of Probability.} Chelsea. — Aksiyomatik olasılığın orijinal metni.

  \item[Lehmann-Casella (1998)] Erich Lehmann, George Casella. \textit{Theory of Point Estimation.} 2nd ed. Springer. — UMVUE, Cramér-Rao, semiparametrik etkinlik.

  \item[Lehmann-Romano (2005)] Erich Lehmann, Joseph Romano. \textit{Testing Statistical Hypotheses.} 3rd ed. Springer. — Neyman-Pearson, UMP, ORT için kapsamlı referans.

  \item[Mood-Graybill-Boes (1974)] A. Mood, F. Graybill, D. Boes. \textit{Introduction to the Theory of Statistics.} 3rd ed. McGraw-Hill.

  \item[Shiryaev Vol.I (1996)] Albert Shiryaev. \textit{Probability.} 2nd ed. Springer. — Rus geleneği; stokastik süreçler giriş.

  \item[Shiryaev Vol.II (2019)] Albert Shiryaev. \textit{Probability-2.} Springer. — Martingaller, ergodik teori.

  \item[Williams (1991)] David Williams. \textit{Probability with Martingales.} Cambridge. — Martingal teorisinin en anlaşılır sunumu; Böl.~\ref{chp:martingaller} için birincil.
\end{description}

\section{İleri Okuma}

\subsection{Olasılık Teorisi}

\begin{description}
  \item[Revuz-Yor (1999)] D. Revuz, M. Yor. \textit{Continuous Martingales and Brownian Motion.} Springer. — Stokastik integral, Itô formülü.
  \item[Kallenberg (2002)] O. Kallenberg. \textit{Foundations of Modern Probability.} 2nd ed. Springer. — Kapsamlı, modern; ileri düzey referans.
  \item[Bass (2011)] Richard Bass. \textit{Stochastic Processes.} Cambridge. — Uygulamalı stokastik süreçler.
\end{description}

\subsection{İstatistik}

\begin{description}
  \item[van der Vaart (1998)] A. W. van der Vaart. \textit{Asymptotic Statistics.} Cambridge. — Donsker, LAN, semiparametrik etkinlik; Böl.~\ref{chp:asimptotik}'ın doğal devamı.
  \item[Wainwright (2019)] M. Wainwright. \textit{High-Dimensional Statistics.} Cambridge. — Büyük boyutlu modern istatistik.
  \item[Efron-Hastie (2016)] B. Efron, T. Hastie. \textit{Computer Age Statistical Inference.} Cambridge. — Bootstrap, lasso, makine öğrenmesi köprüsü.
\end{description}

\subsection{Türkçe Kaynaklar}

\begin{description}
  \item[Akdeniz (2013)] Yukarıda.
  \item[Önder vd. (çeşitli)] Türkiye'deki çeşitli üniversitelerin olasılık-istatistik ders notları (erişim: YÖK Açık Erişim).
\end{description}

\section{Yazılım Kaynakları}

\begin{description}
  \item[R (CRAN)] \texttt{stats}, \texttt{MASS}, \texttt{survival} paketleri; parametrik ve parametresiz testler.
  \item[Python] \texttt{scipy.stats}, \texttt{statsmodels}, \texttt{numpy}; simülasyon ve büyük örneklem doğrulama.
  \item[SageMath] Sembolik hesaplama; olasılık dağılımı türevleri, MÜF hesapları.
\end{description}
